\section*{第 \S\arabic{exercisegroup} 节习题}
\phantomsection
\addcontentsline{toc}{section}{第 \protect\S\arabic{exercisegroup} 节习题}

\begin{enumerate}
\def\labelenumi{\arabic{enumi})}
\tightlist
\item
  建立公式
\end{enumerate}

\[
\tan z = \sum_ {n = 1} ^ {\infty} (- 1) ^ {n - 1} 2 ^ {2 n} (2 ^ {2 n} - 1) b _ {2 n} \frac {z ^ {2 n - 1}}{(2 n) !}
\]

\[
\frac {1}{\sin z} = \frac {1}{z} + \sum_ {n = 1} ^ {\infty} (- 1) ^ {n - 1} 2 (2 ^ {2 n - 1} - 1) b _ {2 n} \frac {z ^ {2 n - 1}}{(2 n) !}
\]

其中右端的级数绝对收敛，第一个对 \(|z| < \frac{\pi}{2}\) ，第二个对 \(|z| < \pi\) （把 \(\tan z\) 与 \(1/\sin 2z\) 表成 \(\cot z\) 与 \(\cot 2z\) 的线性组合）。由此推出数 \(\frac{2^{2n}(2^{2n}-1)}{2n} b_{2n}\) 都是整数。（利用下述引理：若在两个绝对收敛级数 \(\sum_{n=0}^{\infty} \alpha_n \frac{z^n}{n!}\) 、\(\sum_{n=0}^{\infty} \beta_n \frac{z^n}{n!}\) 中系数 \(\alpha_n\) 与 \(\beta_n\) 都是整数，则在写成 \(\sum_{n=0}^{\infty} \gamma_n \frac{z^n}{n!}\) 形式的乘积中，诸 \(\gamma_n\) 都是整数。）

\begin{enumerate}
\def\labelenumi{\arabic{enumi})}
\setcounter{enumi}{1}
\tightlist
\item
  建立公式
\end{enumerate}

\[
(n - 1) \mathrm{B} _ {n} (\mathrm{X}) = n (\mathrm{X} - 1) \mathrm{B} _ {n - 1} (\mathrm{X}) - \sum_ {k = 0} ^ {n} \binom {n} {k} b _ {k}   \mathrm{B} _ {n - k} (\mathrm{X})
\]

（把级数 \(Se^{SX}/(e^{S}-1)\) 对 S 求导）。由此推出公式

\[
(2 n + 1) b _ {2 n} = - \sum_ {k = 1} ^ {n - 1} \binom {2 n} {2 k} b _ {2 k}   b _ {2 n - 2 k}
\]

它是关于伯努利数的。

\begin{enumerate}
\def\labelenumi{\arabic{enumi})}
\setcounter{enumi}{2}
\tightlist
\item
  证明：对每个整数 p \textgreater{} 1 ，
\end{enumerate}

\[
\mathrm{B} _ {n} \left(\frac {x}{p}\right) + \mathrm{B} _ {n} \left(\frac {x + 1}{p}\right) + \dots + \mathrm{B} _ {n} \left(\frac {x + p - 1}{p}\right) = \frac {1}{p ^ {n - 1}} \mathrm{B} _ {n} (\mathrm{X}).
\]

\begin{enumerate}
\def\labelenumi{\arabic{enumi})}
\setcounter{enumi}{3}
\tightlist
\item
  \begin{enumerate}
  \def\labelenumii{\alph{enumii})}
  \tightlist
  \item
    证明关系式
  \end{enumerate}
\end{enumerate}

\[
\mathrm{B} _ {n} (1 - \mathrm{X}) = (- 1) ^ {n} \mathrm{B} _ {n} (\mathrm{X})
\]

（利用事实：\(b_{2n - 1} = 0\) （对 \(n > 1\) ），以及关系式

\[
\mathrm{B} _ {n} (1 - \mathrm{X}) - \mathrm{B} _ {n} (- \mathrm{X}) = (- 1) ^ {n} n \mathrm{X} ^ {n - 1}.
\]

\begin{enumerate}
\def\labelenumi{\alph{enumi})}
\setcounter{enumi}{1}
\tightlist
\item
  证明
\end{enumerate}

\[
\mathrm{B} _ {n} \left(\frac {1}{2}\right) = b _ {n} \left(\frac {1}{2 ^ {n - 1}} - 1\right)
\]

（利用 exerc. 3）。

\begin{enumerate}
\def\labelenumi{\alph{enumi})}
\setcounter{enumi}{2}
\tightlist
\item
  证明：当 n 为偶数时，\(\mathrm{B}_{n}(\mathrm{X})\) 在 R 的区间 {[}0, 1{]} 中有两个根；而当 n 为奇数且 \textgreater{} 1 时，\(\mathrm{B}_{n}(\mathrm{X})\) 在点 0 、\(\frac{1}{2}\) 与 1 处有单根，且在 {[}0, 1{]} 的其他任何点处都不为零（利用 b) 及关系式 \(B_{n}^{\prime}=nB_{n-1}\) ）。
\end{enumerate}

\(d\) ) 由 \(c\) ) 推出：当 \(n\) 为偶数时，\(|\mathrm{B}_n(x)|\) 在区间 {[}0, 1{]} 上的最大值是 \(|b_n|\) ；而当 \(n\) 为奇数时，若 \(a_{n}\) 是 \(|\mathrm{B}_n(x)|\) 在 {[}0, 1{]} 上的最大值，则

\[
\frac {4}{n + 1} \left| b _ {n + 1} \right| \left(1 - \frac {1}{2 ^ {n}}\right) \leqslant a _ {n} \leqslant \frac {1}{2} n \left| b _ {n - 1} \right|
\]

（利用中值定理）。

\begin{enumerate}
\def\labelenumi{\arabic{enumi})}
\setcounter{enumi}{4}
\tightlist
\item
  若置 \(S_{n}(x) = \frac{1}{n + 1} (\mathrm{B}_{n + 1}(x) - \mathrm{B}_{n + 1}(0))\) ，则对每个整数 \(a > 0\) 有
\end{enumerate}

\[
\mathrm{S} _ {n} (a) = 1 ^ {n} + 2 ^ {n} + \dots + (a - 1) ^ {n}.
\]

\begin{enumerate}
\def\labelenumi{\alph{enumi})}
\item
  证明：对每个整数 \(n \geqslant 0\) 和每个整数 \(a > 0\) ，有 \(2\mathrm{S}_{2n + 1}(a) \equiv 0 \pmod{a}\) （考虑和 \(k^{2n + 1} + (a - k)^{2n + 1}\) ）。
\item
  设 \(r\) 和 \(s\) 是任意两个整数，证明
\end{enumerate}

\[
\mathrm{S} _ {n} (r s) \equiv s \mathrm{S} _ {n} (r) + n r \mathrm{S} _ {n - 1} (r) \mathrm{S} _ {1} (s) \pmod {r ^ {2}}.
\]

\begin{enumerate}
\def\labelenumi{\alph{enumi})}
\setcounter{enumi}{2}
\tightlist
\item
  设 p 是素数。证明：若 n 被 p-1 整除，则有 \(S_{n}(p) \equiv -1 \pmod{p}\) ；若 n 不被 p-1 整除，则 \(S_{n}(p) \equiv 0 \pmod{p}\) （若整数 g 不被 p 整除，注意 \(S_{n}(p) \equiv g^{n}S_{n}(p) \pmod{p}\) ）。
\end{enumerate}

\begin{enumerate}
\def\labelenumi{\arabic{enumi})}
\setcounter{enumi}{5}
\tightlist
\item
  \begin{enumerate}
  \def\labelenumii{\alph{enumii})}
  \tightlist
  \item
    设有理数 \(b_{n}\) 已由 VI, p.~275 的公式 (20) 定义，用 \(d_{n}\) 表示 \textgreater{} 0 的分母（即写成既约分数的 \(b_{n}\) 的分母）。证明 \(d_{n}\) 的任何素因子都不能 \(> n + 1\) （利用 VI, p.~276 的递推公式 (23)）。
  \end{enumerate}
\end{enumerate}

\begin{enumerate}
\def\labelenumi{\alph{enumi})}
\setcounter{enumi}{1}
\tightlist
\item
  证明：对每个整数 p \textgreater{} 0 和每个整数 n \textgreater{} 0 ，
\end{enumerate}

\[
\mathrm{S} _ {n} (p) = b _ {n} p + \binom {n} {1} \frac {p}{2} b _ {n - 1} p + \dots + \binom {n} {r} \frac {p ^ {r}}{r + 1} b _ {n - r} p + \dots + \frac {p ^ {n + 1}}{n + 1}.
\]

\begin{enumerate}
\def\labelenumi{\alph{enumi})}
\setcounter{enumi}{2}
\item
  由 b) 出发对 n 作递归，推出对每个素数 p ，把 \(S_{n}(p) - b_{n}p\) 写成既约分数时的分母不被 p 整除（注意 \(p^{r+1}\) 不能整除 \(r + 1\) ）。
\item
  由 c) 推出，数
\end{enumerate}

\[
b _ {n} - \sum_ {p} \frac {\mathrm{S} _ {n} (p)}{p}
\]

（其中 \(p\) 取遍素数集合 \(p \leqslant n + 1\) ，且 \(n\) 为偶数）是整数。由此推出

\[
b _ {2 n} + \sum_ {p} \frac {1}{p}
\]

（其中 \(p\) 取遍满足 \(p - 1\) 整除 \(2n\) 的素数集合）是整数（克劳森–冯·施陶特定理；利用 exerc. 5 c）。

* 7) 我们承认：对每个整数 \(a > 0\) ，整数集合 \(1 + ma\) （ \(m\) 取遍整数集 \(\geqslant 1\) ）中有无穷多个素数（这是狄利克雷关于算术数列的定理的一个特例）。

\begin{enumerate}
\def\labelenumi{\alph{enumi})}
\item
  设 n 是整数 \(\geqslant 1\) ，\(s \geqslant 1\) 是使 \(q = 1 + (2n + 1)!s\) 为素数的整数；证明：若 p 是素数且 p - 1 整除 2nq ，则 p - 1 必整除 2n （相反的情形将有 p - 1 = qd ，其中 d 是整数，从而 p 被 \(d + 1\) 整除）。
\item
  由 a) 推出：对每个整数 n \textgreater{} 0 ，存在无穷多个整数 m \textgreater{} n 使 \(b_{2m} - b_{2n}\) 是整数。*
\end{enumerate}

\begin{enumerate}
\def\labelenumi{\arabic{enumi})}
\setcounter{enumi}{7}
\item
  证明：对每个素数 \(p > 3\) ，把 \(\mathrm{S}_{2n}(p^k) - p^k b_{2n}\) 写成既约分数后，其分子被 \(p^{2k}\) 整除（按 exerc. 6 的方式论证）。
\item
  称有理数 \(r\) 是 \(p\) 进整数（Gen.~Top., III, p.~322, exerc. 23）（\(p\) 为素数），意思是当把 \(r\) 表成既约分数时其分母不被 \(p\) 整除；记 \(r \equiv 0 \pmod{p}\) 以表示 \(r/p\) 是 \(p\) 进整数这一事实，而 \(r \equiv r' \pmod{p}\) 按定义等价于 \(r - r' \equiv 0 \pmod{p}\) 。
\end{enumerate}

\begin{enumerate}
\def\labelenumi{\alph{enumi})}
\tightlist
\item
  设 \(m\) 是有理整数；函数 \(F_{m}(z) = \frac{m}{e^{mz} - 1} - \frac{1}{e^{z} - 1}\) 在 \(|z|\) 充分小时解析，因而可以写成在 0 的一个邻域上收敛的幂级数
\end{enumerate}

\[
\mathrm{F} _ {m} (z) = \sum_ {n = 1} ^ {\infty} (m ^ {n} - 1) \frac {b _ {n}}{n} \frac {z ^ {n - 1}}{(n - 1) !}.
\]

证明在 0 的一个邻域上也有

\[
\mathrm{F} _ {m} (z) = \sum_ {n = 0} ^ {\infty} c _ {n} (e ^ {z} - 1) ^ {n}
\]

其中 \(c_{n}\) 是 \(p\) 进整数；由此推出数 \(a_{n} = (m^{n} - 1)\frac{b_{n}}{n}\) 是 \(p\) 进整数；此外，有 \(a_{n + p - 1} \equiv a_{n} (\mathrm{mod.} p)\) 。

\[
\frac {b _ {n + p - 1}}{n + p - 1} \equiv \frac {b _ {n}}{n} \pmod {p}
\]

\begin{enumerate}
\def\labelenumi{\alph{enumi})}
\setcounter{enumi}{1}
\tightlist
\item
  由 a) 推出：若 \(p - 1\) 不整除 \(n\) ，则 \(b_{n} / n\) 是 \(p\) 进整数，并且有（库默尔同余）。（取 \(m\) 为 mod. \(p\) 类是 \(\mathbf{Z} / p\mathbf{Z}\) 的可逆元素所成乘法群的生成元的整数（Alg., VII. 12）。）
\end{enumerate}

\begin{enumerate}
\def\labelenumi{\arabic{enumi})}
\setcounter{enumi}{9}
\item
  用 exerc. 9 的记号，证明数 \(m^n a_n = m^n (m^n - 1)b_n / n\) 都是整数（把 \(F_m(z)\) 写成 \(e^z\) 的两个多项式之商）。由此推出（对偶数 \(n\) \(\geqslant 2\) ）：分母 \(\delta_n\) （既约形式的 \(b_n / n\) 的）满足：对每个整数 \(m\) （与 \(n\) 互素的）有 \(m^n \equiv 1\) (mod. \(\delta_n\) )。进而，若整数 \(d\) 满足 \(m^n \equiv 1\) (mod. \(d\) )（对每个整数 \(m\) ，与 \(n\) 互素的），则 \(d\) 整除 \(\delta_n\) （利用克劳森–冯·施陶特定理以及 \(\mathbf{Z} / d\mathbf{Z}\) 的乘法群的结构（Alg., VII. 12））。
\item
  \begin{enumerate}
  \def\labelenumii{\alph{enumii})}
  \tightlist
  \item
    设 p 是素数 \(\neq 2\) 。则下列性质等价：
  \end{enumerate}
\end{enumerate}

\(\alpha) b_n / n \not\equiv 0 \pmod{p}\) 对每个偶数 \(n\) （使 \(p - 1\) 不整除 \(n\) 的）。

\(\beta) b_n \not\equiv 0 \pmod{p}\) （对 \(n = 2, 4, 6, \ldots, p - 3\) ）。

（利用 exerc. 9b)。）

\begin{enumerate}
\def\labelenumi{\alph{enumi})}
\setcounter{enumi}{1}
\tightlist
\item
  若 \(p\) 等于 2 或满足 \(a\) ) 中的等价条件，则称 \(p\) 是正则素数；否则称为非正则素数。素数 2, 3, 5, 7, 11 是正则的；691 是非正则的。
\end{enumerate}

设 \(I\) 是非正则素数的一个有限集。设 \(n\) 是整数 \(\geqslant 2\) ，为 \(\prod_{q \in I} (q - 1)\) 的倍数

且使 \(|b_n / n| > 1\) （参看 VI, p.~288，公式 (15)）。设 \(p\) 是既约形式的 \(|b_n / n|\) 的分子的一个素因子。证明 \(p - 1\) 不整除 \(n\) ，因而 \(p\) 是非正则的。由此推出非正则素数的集合是无穷的。

\begin{enumerate}
\def\labelenumi{\arabic{enumi})}
\setcounter{enumi}{11}
\tightlist
\item
  对每个多项式 \(\mathbf{Q}\) （具有实系数且满足 \(\mathrm{Q}(0) = \mathrm{Q}(1)\) 的），用 \(\tilde{\mathbf{Q}}\) 表示 \(\mathbf{R}\) 上满足下列条件的连续函数：\(\tilde{\mathbf{Q}}(x) = \mathbf{Q}(x)\) （对 \(0 \leqslant x \leqslant 1\) ），且 \(\tilde{\mathbf{Q}}(x + 1) = \tilde{\mathbf{Q}}(x)\) （对每个 \(x \in \mathbf{R}\) ）。
\end{enumerate}

\begin{enumerate}
\def\labelenumi{\alph{enumi})}
\item
  对每个整数 \(m \geqslant 2\) ，伯努利多项式 \(B_{m}\) 是唯一的实系数多项式 Q ：次数为 m 、首项系数为 1 ，满足 \(\int_{0}^{1} Q(x) dx = 0\) 、\(Q(1) = Q(0)\) ，且使函数 \(\tilde{Q}\) 在 R 上 m - 2 次可导并有连续的 \((m - 2)^{nd}\) 阶导数。（对 m 作递归论证。）
\item
  证明：对 \(m \geqslant 1\) ，
\end{enumerate}

\[
\int_ {0} ^ {1} \mathrm{B} _ {m} (x) e ^ {- 2 \pi i n x} d x = \left\{ \begin{array}{c c c} 0 & \text {if} & n = 0 \\ - \frac {m !}{(2 \pi i n) ^ {m}} & \text {if} & n \neq 0 \end{array} \right. \text {on} \mathbf {Z}.
\]

* c) 由 b) 推出：对 \(0 \leqslant x \leqslant 1\) 和 \(m \geqslant 2\) ，有

\[
\mathrm{B} _ {m} (x) = - \frac {m !}{(2 \pi i) ^ {m}} \sum_ {n \in \mathbf {Z} - \{0 \}} \frac {e ^ {2 \pi i n x}}{n ^ {m}}
\]

级数正规收敛。*

\begin{enumerate}
\def\labelenumi{\arabic{enumi})}
\setcounter{enumi}{12}
\tightlist
\item
  设 \(f\) 是整数 \(\geqslant 1\) ，\(\chi\) 是定义在 \(\mathbf{Z}\) 上、取复值、且满足 \(\chi(x + f) = \chi(x)\) （对每个 \(x \in \mathbf{Z}\) ）的函数。若置
\end{enumerate}

\[
\hat {\chi} (y) = \frac {1}{f} \sum_ {x = 0} ^ {f - 1} \chi (x) e ^ {- 2 \pi i x y / f} \quad \text { for } y \in \mathbf {Z},
\]

则有

\[
\chi (x) = \sum_ {y = 0} ^ {f - 1} \hat {\chi} (y) e ^ {- 2 \pi i x y / f} \quad \text {   for   } x \in \mathbf {Z}.
\]

由 exerc. 12 推出：对每个整数 \(m \geqslant 2\) ，

\[
\sum_ {n \in \mathbf {Z} \backslash \{0 \}} \frac {\chi (n)}{n ^ {m}} = - \frac {(2 \pi i) ^ {m}}{m !} \sum_ {y = 0} ^ {f - 1} \hat {\chi} (y) \mathrm{B} _ {m} \left(\frac {y}{f}\right).
\]

例如，有

\[
1 - \frac {1}{3 ^ {3}} + \frac {1}{5 ^ {3}} - \dots + \frac {(- 1) ^ {n}}{(2 n + 1) ^ {3}} + \dots = \frac {\pi^ {3}}{3 2}.
\]

\setcounter{exercisegroup}{2}
\refstepcounter{exercisegroup}
